\section{Resource Demo and Case Study}

The demo is designed as a diagnostic-resource check rather than a leaderboard.
It asks whether \tool can ingest real public \sid artifacts from both a
canonical residual-quantization line and a recent recommender-native line, then
produce comparable audit tables with explicit caveats. GRID provides the
open RQ-KMeans-style path~\cite{ju2025grid}; ReSID provides the GAOQ line and
processed Musical artifacts~\cite{liang2026resid}. Sanity tokenizers are kept
as controlled references for metric sensitivity.

\paragraph{Case-study protocol.}
The paper-facing case study uses the same 23,742 Musical items for both rows.
The GRID row is a controlled feature-text export: ReSID processed feature text
is embedded and passed through the official GRID-style residual k-means module.
The ReSID row is a bounded FAMAE-to-GAOQ export on the same item universe. This
protocol is intentionally modest. It gives reviewers a same-item diagnostic
contrast while avoiding the stronger, unsupported claim that the paper
reproduces every training choice of TIGER, GRID, or the full ReSID pipeline.

\begin{table}[t]
\caption{Same-item Musical artifact profiles on 23,742 items. Higher unique
full codes, D3 L1 co-occurrence prefix recall, and D4 tail unique-SID ratio are
better; lower D2 full-collision rate is better. D5a prefixes are level-wise
prefix counts, not latency. This is not a downstream leaderboard: GRID
feature-text uses ReSID processed feature text, and ReSID is a bounded GAOQ
export.}
\label{tab:musical-diagnostic}
\small
\setlength{\tabcolsep}{2.0pt}
\begin{tabular}{@{}lrrrrl@{}}
\toprule
System & Unique & D2 full & D3 L1 & D4 tail & D5a prefixes \\
\midrule
GRID feature-text & 3,749 & 0.9769 & 0.0552 & 0.3695 & 64/3440/3749 \\
ReSID bounded & 23,742 & 0.0000 & 0.1535 & 1.0000 & 32/1280/23742 \\
\bottomrule
\end{tabular}
\end{table}

Table~\ref{tab:musical-diagnostic} compares two exports on the same 23,742
Musical items. The GRID row is intentionally labeled as feature-text: it uses
processed feature text from the ReSID artifact path and the official
MiniBatchKMeans-style residual quantization module, not a faithful raw-text
TIGER reproduction. Under that controlled setup, the row exposes high collision
pressure and limited tail capacity. The ReSID/GAOQ row exports one full code
per item in this bounded run, yielding zero full-code collisions and higher
D3-L1 collaborative prefix recovery. The D5a prefix column reports level-wise
prefix counts, a structural proxy for trie expansion rather than real serving
latency.

This contrast is useful precisely because the table is not a downstream
ranking claim. It shows that the same audit interface can surface different
artifact mechanisms: compact residual-k-means codes can collapse many items
into shared full \sid{}s; recommender-native GAOQ can avoid full collisions in
the exported mapping; and the tail-capacity column makes allocation pressure
visible rather than inferred from item popularity. Additional artifact tables
cover sanity controls, GRID scale checks, DACT churn, and a MovieLens schema
smoke, but those are placed in the repository rather than the four-page PDF.

\paragraph{What the extra artifact tables add.}
The repository tables are not backup evidence for hidden claims; they serve
specific resource checks that do not need space in the PDF. Sanity controls
show that D1--D5a react to intentionally extreme tokenizers rather than
collapsing to constant values. GRID scale checks show that the adapter and
metric runner handle larger item counts than the Musical case study. The DACT
table demonstrates that the optional D6 churn interface can compare two SID
refreshes. The MovieLens smoke test verifies that the schema is not tied to
Amazon metadata. These tables support the artifact's usability claim, whereas
Table~\ref{tab:musical-diagnostic} remains the only method-facing diagnostic
case study in the main paper.
